\section{The MMVSD algorithm}\label{sec:mmvsdalg}

\subsection{Frames of several lengths}\label{sec:frames}

The modulated signal is built from short pieces of noise, \emph{frames}, each with its own spectrum
$m(f,t_\ell)$. A frame of length $L$ resolves frequency to $1/L$ and time to about $L$: long frames
are needed at low frequencies (an oscillator needs several cycles in a frame), short ones to follow
short bursts. The program therefore uses frames whose length grows towards low frequencies.

\emph{Requirements.} Node $i$ needs a hop (the spacing of consecutive frames, half the frame length)
of at most a third of its burst duration, and frames of at least 8 cycles, i.e.\ a hop of at least
$4/f_{\mathrm{n},i}$. The hop it gets is the larger of the two, $h_i = \max(T_{c,i}/3,\
4/f_{\mathrm{n},i})$, made non-increasing towards high frequencies (a cumulative minimum from the
lowest node up).

\emph{Groups.} The finest hop $h_0 = \min_ih_i$ (an even number of samples, at least 16) defines the
fine time grid. Every node is assigned to a group $g\in\{0,1,\dots\}$ with the hop $h_02^{g}$: the
longest power-of-two multiple of $h_0$ that does not exceed its need $h_i$, but long enough for the
cycle condition. At most six groups are used, the coarsest with at least 16 frames in the test, and the
test length is rounded to a whole number of coarse frames.

\emph{Crossovers.} Each group synthesises the nodes of its frequency range. Two adjacent groups are
joined at a boundary frequency $f_b$ (geometric mean of their edge nodes) with power-complementary
weights over 1/6 octave: with
\begin{equation*}
  t = \min\Bigl(1,\ \max\Bigl(0,\ 6\log_2\frac{f}{f_b} + \frac12\Bigr)\Bigr)
\end{equation*}
the lower group receives $1-\sin^2(\pi t/2)$ and the upper one $\sin^2(\pi t/2)$ of the power. Applied
recursively from the lowest boundary up, this defines weights $X_g(f)\ge0$ with $\sum_gX_g(f) = 1$
at every frequency.

\subsection{Synthesis of one group}

Let the group have the hop $h$ (in samples) and the frame length $L = 2h$. Frame $\ell$ covers the
samples $(\ell-1)h,\dots,(\ell+1)h-1$ (circularly) and is multiplied by the \emph{square-root Hann}
window $w(n) = \sin\bigl(\pi(n+\tfrac12)/L\bigr)$, $n = 0,\dots,L-1$. At any sample two frames overlap,
at positions $n$ and $n+h$ of their windows, and
\begin{equation}
  w(n)^2 + w(n+h)^2 = \sin^2\!\Bigl(\frac{\pi(n+\frac12)}{L}\Bigr) +
  \sin^2\!\Bigl(\frac{\pi(n+\frac12)}{L}+\frac{\pi}{2}\Bigr)
  = \sin^2(\cdot)+\cos^2(\cdot) = 1,
\end{equation}
because $\pi h/L = \pi/2$. The squared windows of overlapping frames therefore add up to one.

Each frame is a periodic noise built from its discrete Fourier coefficients on the lines $f_q =
q\,f_s/L$, $q = 0,\dots,h$:
\begin{equation}
  C_{\ell q} = \sqrt{\tfrac{Lf_s}{2}}\;\sqrt{X_g(f_q)}\;\sqrt{m_\ell(f_q)}\;c_{\ell q},
\end{equation}
where $c_{\ell q}$ are random coefficients with $\E|c_{\ell q}|^2=1$ (Section~\ref{sec:framenoise}) and
$m_\ell(f_q)$ is the field of the group (the fine field averaged over the $2^g$ fine frames of the
frame), interpolated linearly in $\log_2 f$ between the nodes. The inverse transform of length $L$
turns the coefficient $C$ of line $q$ into a cosine of amplitude $2|C|/L$, i.e.\ of power
$\tfrac12(2|C|/L)^2 = 2|C|^2/L^2$. With the factor $\sqrt{Lf_s/2}$ this power is $(f_s/L)\,X_gm_\ell|c|^2$,
and dividing by the line spacing $f_s/L$ gives the PSD $X_g(f)\,m_\ell(f)$ (on average over $c$):
with $m=1$ and $X_g=1$ the frame is white noise of unit one-sided PSD.

The frames of a group are independent and have zero mean, so at a given time the variances of the
two overlapping (windowed) frames add. Summing over the groups (also independent),
\begin{result}[Local spectrum of the synthesised noise]
\begin{equation}\label{eq:meff}
  m_{\mathrm{eff}}(f,t) = \sum_g X_g(f)\sum_\ell w_{g\ell}(t)^2\,m_{g\ell}(f),
\end{equation}
\end{result}
which is 1 for $m=1$ (since $\sum_\ell w^2 = 1$ and $\sum_g X_g = 1$) and has time average 1 when every
row of $m$ has mean 1. The model evaluates~\eqref{eq:meff} on the fine grid with two points per fine
hop ($\delta = h_0\Delta t/2$): between two frame centres the weights are $\cos^2(\pi u/2)$ and
$\sin^2(\pi u/2)$, $u\in(0,1)$ being the position between them, exactly the squared windows.

\emph{Shaping.} The sum of the groups, $v(t)$, is a noise whose local spectrum is
$m_{\mathrm{eff}}(f,t)$. Its discrete Fourier transform over the whole test, multiplied by
$\sqrt{P(f)} = \sqrt{\Pi(f)\,p(f)}$ on every line, gives the test signal: a zero-phase filter that
turns the unit spectrum into $P(f)$. The local spectrum of the result is $P(f)\,m_{\mathrm{eff}}(f,t)$,
which is the model~\eqref{eq:evol} with $m = m_{\mathrm{eff}}$.

\subsection{The frame noise}\label{sec:framenoise}

Two kinds of coefficients $c_{\ell q}$ are available.
\begin{itemize}
\item \emph{Gaussian frames}: complex normal coefficients (independent real and imaginary parts of
variance $\tfrac12$). Every frame is then a Gaussian noise, and so is the sum of frames and its
filtered version at every instant, for any frame length: the signal is conditionally Gaussian, as
the model assumes (Section~\ref{sec:multisine}).
\item \emph{Multisine frames}: coefficients of modulus 1 with random phases. Each frame has exactly
the band power of its spectrum, so the realisation scatters less around the model; but when a frame
contains only a few lines in the band of an oscillator ($L\,\fn/Q$ small, typically at low frequency
with short frames), the response within the frame is a sum of a few sinusoids whose amplitudes are
fixed: its peaks are lighter-tailed than Gaussian, and its damage is on average 20 to 40\,\% lower than the
Gaussian value at the same variance. The closed loop compensates with a higher level and a deeper
modulation.
\end{itemize}
The program uses Gaussian frames (\code{frame\_noise = "gaussian"}), so that the realisations
follow the statistics that the model assumes. On measured transport records (not distributed with
the package) they gave lower maximum FDS errors and lower $\kappa_z$ errors after the closed loop than
multisine frames; on synthetic records the two were within the scatter between seeds. The scatter
itself is a property of the target: with $b = 8$ the damage of a strongly modulated response rests on
its few largest bursts, and the damage of a single oscillator typically varies by a factor of about
1.5 between two noises of the same field (standard deviation of about 0.5 in $\ln D$ for a field of
depth $e = 3$ in a 60\,s test). This is the main reason for the closed loop of
Section~\ref{sec:mmvsdloop}.

\subsection{Design}\label{sec:design}

The design finds the nodal PSD factors $p_j$ and modulation depths $e_j$ for which the model meets the
targets, before any signal is synthesised. The unknowns are taken in logarithms, $y_j = \ln p_j$ and
$\chi_j = \ln e_j$, which keeps them positive and makes relative changes uniform. With the field
fixed, the model gives the residuals
\begin{equation}
  r_{D,i}(y) = \ln\frac{D_i(y)}{D_{t,i}}, \qquad r_{\kappa,i}(\chi) = \kappa_i(\chi) - \kappa_{t,i}.
\end{equation}
Since the kurtosis depends mainly on the depth and the damage on both, the two are solved
alternately, three times: first $\chi$ for $r_\kappa = 0$ with $y$ fixed, then $y$ for $r_D=0$ with
$\chi$ fixed.

Each problem is a regularised least-squares problem solved by the Levenberg--Marquardt method. For the
level, for example,
\begin{equation}
  \min_y\ \|r_D(y)\|^2 + \mu\,\|L(y - y_0)\|^2 ,
\end{equation}
where $L$ is the second-difference matrix, $(Ly)_i = y_i - 2y_{i+1} + y_{i+2}$, $\mu$ the weight of
the curvature penalty ($0.05$ for the level, $0.5$ for the depth) and $y_0$ the preferred solution
($y_0=0$ with the shape prior: no correction of the reference shape unless the targets require it).
Linearising $r_D(y+\delta) \approx r_D + J\delta$, with the Jacobian $J_{ij} = \partial r_{D,i}/\partial
y_j$, and setting the gradient of the objective to zero gives the step
\begin{equation}
  \bigl(J^{\top}J + \mu L^{\top}L + \lambda\,\mathrm{diag}(J^{\top}J+\mu L^{\top}L)\bigr)\,\delta
  = -\bigl(J^{\top}r_D + \mu L^{\top}L(y-y_0)\bigr).
\end{equation}
The damping $\lambda$ is divided by 3 after a step that lowers the objective and multiplied by 4 after
one that does not (then the step is retried). The depths are limited to $10^{-3}\le e_j\le e_{\max}$.
The level design starts from the median level of the stationary Gaussian solution of
Section~\ref{sec:inversion}, the depth design from $e_j = \kappa_{t,j}-3$.

\subsection{Exact Jacobians}\label{sec:jacobians}

The model~\eqref{eq:S}--\eqref{eq:Dmodel} is a composition of linear operators and pointwise
functions, so its derivatives can be computed exactly and cheaply.

\emph{Damage with respect to the level.} On the model grid of $N_t$ points, $D_i =
\frac{K^b}{C}\frac{T}{N_t}\sum_l d\bigl(S_i(t_l)\bigr)$. Let $g_{k,i}(t) = \partial d/\partial
S_k$ evaluated at $S_i(t)$. By~\eqref{eq:S}, $\partial S_{k,i}(t)/\partial y_j = W^{(k)}_{ij}
F_i[Y_j](t)$ with $Y_j(t) = p_jm_j(t)$ (because $\partial p_j/\partial y_j = p_j$). By the chain rule,
\begin{equation}
  \frac{\partial D_i}{\partial y_j} = \frac{K^b}{C}\frac{T}{N_t}\sum_k W^{(k)}_{ij}\sum_l g_{k,i}(t_l)\,F_i[Y_j](t_l).
\end{equation}
The inner sum is a scalar product between $g_{k,i}$ and a filtered signal. For any linear filter
$F$, $\sum_l g(t_l)\,F[Y](t_l) = \sum_l F^{\top}[g](t_l)\,Y(t_l)$, where the \emph{adjoint} $F^{\top}$ is
the same filter run backwards in time (in the Fourier domain: multiplication by the complex conjugate
of its response). Hence
\begin{result}[Damage Jacobian]
\begin{equation}
  \frac{\partial\ln D_i}{\partial y_j} = \frac{1}{D_i}\frac{K^b}{C}\frac{T}{N_t}
  \sum_k W^{(k)}_{ij}\sum_l F_i^{\top}[g_{k,i}](t_l)\;p_j\,m_j(t_l).
\end{equation}
\end{result}
The $4n$ adjoint filterings and one matrix product with the field give all $n^2$ entries at the cost
of a few model evaluations. The derivatives $g_{k,i}$ of the closed-form rate $d$ are computed by
one-sided differences with a relative step of $10^{-5}$.

\emph{Kurtosis with respect to the depth.} With $A_i = \tavg{s_i^2}$ and $B_i = \tavg{s_i}$,
$\kappa_i = 3A_i/B_i^2$. A change $\delta m_j(t)$ of the field row $j$ changes $s_i$ by $\delta s_i
= F_i[W^{(0)}_{ij}p_j\,\delta m_j]$. The mean $B_i$ does not change, because the field keeps
$\tavg{m_j} = 1$ and therefore $\tavg{\delta m_j} = 0$. The mean square changes by
$\delta A_i = 2\tavg{s_i\,\delta s_i} = 2W^{(0)}_{ij}p_j\tavg{F_i^{\top}[s_i]\,\delta m_j}$, so
\begin{equation}
  \frac{\partial\kappa_i}{\partial\chi_j} = \frac{6\,W^{(0)}_{ij}\,p_j}{B_i^2}\,
  \Bigl\langle F_i^{\top}[s_i]\;\frac{\partial m_j}{\partial\chi_j}\Bigr\rangle ,
\end{equation}
where $\partial m_j/\partial\chi_j$ is obtained by rebuilding the field with $e_j$ multiplied by
$e^{0.05}$ (the field depends smoothly on $e_j$). The same steps give the derivative of $\ln D_i$ with
respect to the depth, used by the closed loop.

\emph{Kurtosis with respect to the level.} Here $B_i$ changes too: $\partial B_i/\partial y_j =
W^{(0)}_{ij}p_j\tavg{m_j}$ and $\partial A_i/\partial y_j = 2W^{(0)}_{ij}p_j\tavg{F_i^\top[s_i]\,m_j}$,
and $\partial\kappa_i = 3\bigl(\partial A_i/B_i^2 - 2A_i\,\partial B_i/B_i^3\bigr)$.

\subsection{The closed loop on the realisation}\label{sec:mmvsdloop}

The design is synthesised once with frozen frame coefficients; from then on the signal is a smooth
function of $(p,e)$. Every iteration measures, on the realisation, the rainflow FDS $D_{m}$ and the
kurtosis of the responses $\kappa_{m}$ (both in one pass over the oscillators) and forms
\begin{equation}
  r_{D,i} = \ln\frac{D_{m,i}}{D_{t,i}},\qquad r_{\kappa,i} = w_\kappa\,\bigl(\bar\kappa_{m,i}-\kappa_{t,i}\bigr),
\end{equation}
with $w_\kappa = 0.15$ and $\bar\kappa_m$ the measured kurtosis averaged over three adjacent nodes.
The quality of a realisation is the score $\tfrac12\max_i|r_{D,i}| + \tfrac12q_{90}(|r_D|) +
\mathrm{mean}_i|r_{\kappa,i}|$, and the best realisation found is kept. Four kinds of steps are used.

\emph{Early re-draws.} The first realisation of the design can miss the target damage of one
oscillator by a large factor (several hundred per cent) while the median error is a few tens of per
cent: with $b$ between 5 and 10, one unusually large burst of the noise within the band of that
oscillator carries most of its damage. A correction of the level would then dig a notch into $P$ at
that node to compensate one rare event of the noise. Before any correction, therefore, while the worst
oscillator misses its damage by more than 50\,\%, the response $z$ of that oscillator is computed on
the realisation, the blocks of its group where $\sum|z|^b$ is largest are found, and the frames that
build them (the two largest blocks and their neighbours) receive new random coefficients on the lines
within $\pm1/4$ octave of the oscillator. The new draw is kept only if it lowers both the score and
the error of that oscillator; otherwise the old coefficients are restored. At most six draws are made,
and the phase ends after two useless draws in a row, because the remaining error is then a property
of the design and not of the noise.

\emph{Coupled Gauss--Newton steps} (first phase of the corrections). The model Jacobians carry the couplings between
oscillators (shared bands, the modulation of a neighbouring band, the envelope law), the measured
residuals what the model neglects. First the depth moves to correct the kurtosis,
\begin{equation}
  \delta\chi = -\bigl(J_{\kappa\chi}^{\top}J_{\kappa\chi} + \mu'L^{\top}L + \lambda\,\mathrm{diag}\bigr)^{-1}
  \bigl(J_{\kappa\chi}^{\top}r_\kappa + \mu'L^{\top}L\chi\bigr),
\end{equation}
then the level moves to correct the damage, including the change that the model predicts for the new
depth, $r_D + J_{D\chi}\delta\chi$:
\begin{equation}
  \delta y = -\bigl(J_{Dy}^{\top}J_{Dy} + \mu''L^{\top}L + \lambda\,\mathrm{diag}\bigr)^{-1}
  \bigl(J_{Dy}^{\top}(r_D + J_{D\chi}\delta\chi) + \mu''L^{\top}L(y-y_0)\bigr),
\end{equation}
with weak curvature weights ($\mu'' = 10^{-4}$, $\mu' = 0.05$) so that node-wise corrections remain
possible, and steps limited to $\pm0.5$ (depth) and $\pm0.3$ (level) in logarithm. After two steps
that do not improve the score, the loop switches to polishing.

\emph{Polishing.} Only the level moves, by the damped least-squares step
\begin{equation}
  \delta y = -w\,\bigl(J_{Dy}^{\top}J_{Dy} + \mu_p\,\mathrm{diag}(J_{Dy}^{\top}J_{Dy})\bigr)^{-1}J_{Dy}^{\top}r_D ,
  \qquad \mu_p = 0.5 .
\end{equation}
If the Jacobian were diagonal with the entries $b/2$ (each oscillator responding only to its own
node, $D\propto p^{b/2}$), this would be $\delta y_i = -w\,r_{D,i}/\bigl((b/2)(1+\mu_p)\bigr)$, i.e.\ the
node-wise step~\eqref{eq:gain} of version~2.0 ($\ln g_i^2 = (2w/b)\ln(D_t/D)$) reduced by $1+\mu_p$.
With the full Jacobian, the correction of an oscillator whose damage comes from a neighbouring band
goes to that band instead of driving its own node up or down. The relaxation $w$ adapts as in
Section~\ref{sec:closedloop}.

\emph{Re-draws.} The largest residual errors are usually a property of the realisation: the damage of
an oscillator rests on a few large cycles, and the level cannot fix it without disturbing its
neighbours. When the polishing stalls, the frames that carry most of the model damage of the worst
oscillator (the two largest and their neighbours, in its group) receive new random coefficients on
the lines within $\pm1/4$ octave of it, and a short branch of four polishing steps follows. The branch is
kept only if it beats the best realisation so far; otherwise the old coefficients are restored. At most
eight re-draws are tried.

\emph{Stopping.} The loop stops when the maximum FDS error is below the tolerance and the median
kurtosis error below 0.3, when the iteration budget is used, or when the score has not improved by
1\,\% for ten iterations.

\emph{A second strategy.} The coupled loop corrects kurtosis and damage together and usually
converges in a few tens of iterations. Where the damage of a few oscillators rests on a few bursts it
can stall at a compromise between neighbouring oscillators, because the model Jacobian then
describes the realisation poorly. On request (option \code{fallback}), when its maximum FDS error ends
above twice the tolerance, the program runs a second closed loop from the same design and the same
noise with node-wise steps only (the step of Section~\ref{sec:closedloop}, the depth kept at the
design), which progresses more slowly but steadily, and keeps the realisation with the better score.
The option is off by default: on measured transport records the coupled loop was better in every
case, and the second loop only doubled the computing time. The program warns when the final level differs from the design by more than
a factor of 8 over a band (a systematic drift, typically at oscillators whose damage is produced by
another band), and when the target kurtosis could not be reached with the maximum depth.
